\chapter{Input Generation}
\label{ch:input_gen}


% ======================================================================
\section{Interatomic Force Constants}
\label{sec:ifc}
% ======================================================================


% ----------------------------------------------------------------------
\subsection{Finite-displacement FCs}
\label{sec:finite_displacement}
% ----------------------------------------------------------------------

The FCs can be generated based on a finite-displacement workflow
using \texttt{phono3py}, QE, and the \texttt{symfc} fitting backend.  In this
approach, displaced supercells are generated explicitly, DFT forces are
computed for each displaced configuration, and both FC2 and FC3 are
reconstructed from the resulting force data.

The following steps are performed:
\begin{enumerate}
  \item \textbf{Structure setup.}
        The crystal structure is loaded and a phonopy object is
        created. The unit cell is loaded via \\
        \texttt{structure.py:load\_structure}, and a
        \texttt{Phonopy} object with a supercell is
        created by
        \texttt{structure.py:create\_phonopy\_from\_config}, which
        generates symmetry-inequivalent displacements.
  \item \textbf{DFT force calculations.}
        For each displaced supercell, a QE \texttt{pw.x}
        input is written and executed.  A QE \texttt{pw.x} SCF input
        file is written by\\
        \texttt{qe\_interface.py:write\_qe\_scf\_input}.  Existing
        converged outputs are detected and reused.  Forces parsed from
        QE output are converted from Ry/bohr to eV/\AA.

  \item \textbf{Force-constant reconstruction.}
        The set of forces is passed to phono3py, which reconstructs FC3
        and FC2 using the symmetry-adapted \texttt{symfc} backend, both
        FC2 and FC3 are written to a common \texttt{fc3.hdf5} file
        containing dense arrays in eV/\AA$^2$ and eV/\AA$^3$,
        respectively.
\end{enumerate}

For the Si FCC primitive cell with a $2\times2\times2$ supercell
(16 atoms), phono3py generates 114 symmetry-inequivalent displacements
for a pair cutoff of 5.5\,\AA.  On the mont-fort workstations, each QE SCF
calculation takes roughly 30\,s  with
\texttt{ecutwfc}=60\,Ry and a $2\times2\times2$ $k$-point mesh.

% ----------------------------------------------------------------------
\subsection{DFPT FCs}
\label{sec:dfpt}
% ----------------------------------------------------------------------

Third order FCs can be generated efficiently directly from density-functional
perturbation theory (DFPT) due to the $2n+1$ theorem.  One possibility is to obtain FC2 from QE's linear-response phonon code \texttt{ph.x} and compute FC3 with the D3Q plugin,
which implements third-order DFPT.  In contrast to the
finite-displacement method, the DFPT route works directly on the unit
cell and avoids supercell calculations. However, QE does only supports a restricted set of
norm-conserving potentials for third-order DFPT calculations, limiting the achievable accuracy.

The DFPT workflow does the following:
\begin{enumerate}
  \item \textbf{SCF ground state} (\texttt{pw.x}).
        A self-consistent calculation on the unit cell with a dense
        $k$-mesh (configured via \texttt{dfpt.kpoints}).

  \item \textbf{Phonon calculation} (\texttt{ph.x}).
        DFPT phonon calculation on a uniform $\vc{q}$-grid
        (\texttt{dfpt.q\_mesh}) with
        \texttt{ldisp\,=\,.true.} and \texttt{lqdir\,=\,.true.}.  This
        produces the dynamical matrices $D(\vc{q})$ and saves the
        first-order perturbation wavefunctions needed by D3Q.

  \item \textbf{FC2 Fourier transform} (\texttt{q2r.x}).
        The dynamical matrices are Fourier-transformed to real-space FC2
        via QE's \texttt{q2r.x}, producing the \texttt{flfrc} file.
        The acoustic sum rule is enforced
        (\texttt{zasr\,=\,'crystal'}).

  \item \textbf{Third-order DFPT} (\texttt{d3q.x}).
        For each $(\vc{q}_1,\vc{q}_2)$ pair on the grid, \texttt{d3q.x}
        computes the third-order dynamical tensor
        $D_3(\vc{q}_1,\vc{q}_2,-\vc{q}_1-\vc{q}_2)$.  The
        $\vc{q}$-points are specified in QE Cartesian coordinates
        ($2\pi/a_\text{lat}$ units), the conversion from crystal
        reciprocal coordinates is
        \begin{equation}\label{eq:q_conversion}
          \vc{q}_\text{cart}
          = q_1 \hat{\vc{b}}_1 + q_2 \hat{\vc{b}}_2 + q_3 \hat{\vc{b}}_3,
          \qquad
          \hat{B} = (\mathcal{A} / |\vc{a}_1|)^{-\top},
        \end{equation}
        where $\mathcal{A}$ is the matrix of lattice vectors (rows, in
        \AA; distinct from the \AA-to-metre factor $A$ of
        \cref{sec:phonopy_units}) and
        $\hat{\vc{b}}_i$ are the rows of $\hat{B}$.

        For a $2\times2\times2$ grid there are $8^2 = 64$ pairs.
        Crystal symmetry can reduce this (e.g.\ to ${\sim}10$--15 for
        FCC Si with $O_h$ symmetry). The current implementation
        generates all pairs and relies on \texttt{d3q.x}'s internal
        symmetry handling.

  \item \textbf{FC3 Fourier transform} (\texttt{d3\_qq2rr.x}).
        The $(\vc{q}_1,\vc{q}_2)$-space third-order tensors are
        Fourier-transformed to real-space FC3
        $\Phi_3(\vc{R}_1,\vc{R}_2)$, producing the \texttt{mat3R} file.

  \item \textbf{Acoustic sum rule} (\texttt{d3\_asr.x}).
        Enforces translational invariance on the real-space FC3,
        $\sum_\kappa \Phi_3(\kappa,\kappa',\kappa'') = 0$ for each
        pair $(\kappa',\kappa'')$.  The ASR type is configurable
        (\texttt{asr}: \texttt{simple}, \texttt{crystal}, or
        \texttt{no}).

  \item \textbf{Sparsification} (\texttt{d3\_sparse.x}).
        Zeroes FC3 entries with magnitude below a threshold
        (\texttt{sparse\_thr}, in Ry/bohr$^3$), reducing storage
        and downstream computation cost.
\end{enumerate}

\paragraph{Output conversion.}
The \texttt{reap} step parses both the \texttt{q2r.x} FC2 output
(\texttt{fc2.dat}, in Ry/bohr$^2$) and the D3Q \texttt{mat3R} file (in
Ry/bohr$^3$) and converts to eV/\AA$^2$ and eV/\AA$^3$:
\begin{align}
  \Phi_2\;[\mathrm{eV}/\text{\AA}^2]
    &= \Phi_2\;[\mathrm{Ry/bohr^2}]
       \times \frac{E_h/2}{a_0^2},
    &&\approx 48.59, \label{eq:fc2_conv} \\
  \Phi_3\;[\mathrm{eV}/\text{\AA}^3]
    &= \Phi_3\;[\mathrm{Ry/bohr^3}]
       \times \frac{E_h/2}{a_0^3},
    &&\approx 91.82. \label{eq:fc3_conv}
\end{align}
The real-space lattice vectors $\vc{R}$ in the D3Q output (integer
crystal coordinates) are mapped to supercell atom indices
$j = (l_1 n_2 n_3 + l_2 n_3 + l_3)\, n_\text{at} + \kappa$ with
$l_i = R_i \bmod n_i$, matching phonopy's supercell convention.  FC2
is expanded to the full $(N_\text{super})^2$ array via translational
invariance.  Both arrays are saved to \texttt{fc3.hdf5} in the same
format as the finite-displacement method.

The DFPT-specific parameters are in the \texttt{dfpt} section of the
YAML config:
\begin{center}
\begin{tabular}{lll}
  \toprule
  Key & Default & Description \\
  \midrule
  \texttt{q\_mesh}   & $[2,2,2]$ & $\vc{q}$-grid for \texttt{ph.x} and \texttt{d3q.x} \\
  \texttt{kpoints}   & $[8,8,8]$ & SCF $k$-mesh for the unit cell \\
  \texttt{tr2\_ph}   & $10^{-14}$& \texttt{ph.x} self-consistency threshold \\
  \texttt{ph\_command}       & \texttt{ph.x}          & \texttt{ph.x} executable \\
  \texttt{d3q\_command}      & \texttt{d3q.x}         & \texttt{d3q.x} executable \\
  \texttt{d3\_qq2rr\_command}& \texttt{d3\_qq2rr.x}   & Fourier transform executable \\
  \texttt{d3\_asr\_command}  & \texttt{d3\_asr.x}     & Acoustic sum rule executable \\
  \texttt{d3\_sparse\_command}& \texttt{d3\_sparse.x} & Sparsification executable \\
  \texttt{asr}               & \texttt{simple}        & ASR type: \texttt{simple}, \texttt{crystal}, or \texttt{no} \\
  \texttt{sparse\_thr}       & $10^{-5}$              & FC3 threshold (Ry/bohr$^3$); \texttt{null} to skip \\
  \texttt{q2r\_command}      & \texttt{q2r.x}         & \texttt{q2r.x} executable \\
  \texttt{work\_dir}         & \texttt{./dfpt}        & Working directory \\
  \texttt{ph\_timeout}       & 7200\,s                & Timeout per \texttt{ph.x} run \\
  \texttt{d3q\_timeout}      & 14400\,s               & Timeout per \texttt{d3q.x} triplet \\
  \bottomrule
\end{tabular}
\end{center}

% ----------------------------------------------------------------------
\subsection{Mass-weighting and unit conversion}
\label{sec:ifc_mass_weighting}
% ----------------------------------------------------------------------

The harmonic transport problem is formulated in terms of mass-weighted
FC2 blocks.  Given a harmonic IFC block
$\Phi_{\alpha\beta}(0\kappa,l'\kappa')$, the corresponding mass-weighted
quantity is
\begin{equation}
  H_{\alpha\beta}(0\kappa,l'\kappa')
  = \frac{\Phi_{\alpha\beta}(0\kappa,l'\kappa')}
         {\sqrt{M_\kappa M_{\kappa'}}},
\end{equation}
with masses in amu.  In the common phonopy unit system,
$H_{\alpha\beta}$ therefore has units
eV/(amu$\cdot$\AA$^2$), which are later converted to
$(\mathrm{rad/s})^2$ using the factor defined in
\cref{sec:solver_units}.

For FC3, the analogous mass-weighted tensor is
\begin{equation}\label{eq:fc3_mw}
  \Psi_{\alpha\beta\gamma}(\kappa\kappa'\kappa'')
  = \frac{\Phi^{(3)}_{\alpha\beta\gamma}(\kappa\kappa'\kappa'')}
         {\sqrt{M_\kappa M_{\kappa'} M_{\kappa''}}},
\end{equation}
with units
\begin{equation}
  [\Psi]
  = \frac{\mathrm{eV}}
         {\text{\AA}^3\cdot\mathrm{amu}^{3/2}}.
\end{equation}

For DFPT-generated FCs, the real-space force constants are parsed in
atomic units and converted with the same factors as in
\cref{eq:fc2_conv,eq:fc3_conv}.

Finally, the FC3 tensor can be converted to SI units for use in the
anharmonic self-energy.  Using the harmonic conversion factor
$C = e/(m_u A^2)$, the cubic conversion becomes
\begin{equation}\label{eq:conversion_fc3_ifc}
  \texttt{CONVERSION\_FC3}
  = \frac{C}{A\sqrt{m_u}}
  = \frac{e}{m_u^{3/2} A^3}.
\end{equation}
This converts
eV/(\AA$^3\cdot$amu$^{3/2}$) into the corresponding SI units required
by the phonon--phonon self-energy expressions.

% ----------------------------------------------------------------------
\subsection{Validation of the extracted harmonic model}
\label{sec:extraction_validation}
% ----------------------------------------------------------------------

Before using the extracted harmonic blocks in transport calculations, it
is important to verify that the block representation is internally
consistent and reproduces the original phonopy dynamical matrix.  The
present validation procedures are implemented in
\texttt{validation.py}.

\paragraph{Acoustic sum rule.}

At the $\Gamma$ point, the lattice Fourier series reduces to
\begin{equation}
  D^{(\mathrm{B})}(\vc{0}) = \sum_{\vc{R}} H(\vc{R}).
\end{equation}
Translational invariance requires that the three acoustic modes vanish
at $\Gamma$.  The validation therefore forms
\begin{equation}
  D_\Gamma = \sum_{\vc{R}} H(\vc{R}),
\end{equation}
computes its eigenvalues $\lambda_n$, and reports the corresponding
frequencies
\begin{equation}
  \omega_n = \mathrm{sign}(\lambda_n)\sqrt{|\lambda_n|}.
\end{equation}
The three smallest $|\omega_n|$ should be close to zero.  In addition,
the matrix should be symmetric to numerical precision, so that
\begin{equation}
  \max |D_\Gamma - D_\Gamma^T|
\end{equation}
is approximately zero.

\paragraph{Block symmetry.}

The real-space harmonic blocks must satisfy the Hermiticity relation
\begin{equation}\label{eq:block_sym}
  H(\vc{R})^T = H(-\vc{R})
  \qquad \forall\;\vc{R}.
\end{equation}
The validation iterates over all block pairs
$(\vc{R},-\vc{R})$ and reports the maximum deviation
\begin{equation}
  \max_{\vc{R}}
  \|H(\vc{R})^T - H(-\vc{R})\|_\infty.
\end{equation}

\paragraph{Band-structure comparison.}

A more stringent validation compares phonon frequencies reconstructed
from the extracted blocks with those obtained directly from phonopy.
Along a chosen $\vc{q}$-path, two spectra are computed:
\begin{enumerate}
  \item the phonopy reference spectrum, obtained from
        \begin{equation}
          D_\mathrm{ph}(\vc{q})
          = P(\vc{q})\,D^{(\mathrm{A})}(\vc{q})\,P^\dagger(\vc{q}),
        \end{equation}
        converted to $(\mathrm{rad/s})^2$ and diagonalized;

  \item the reconstructed spectrum, obtained from
        \begin{equation}
          D_\mathrm{rec}(\vc{q})
          = \sum_{\vc{R}} H(\vc{R})\,
            e^{2\pi i\,\vc{q}\cdot\vc{R}},
        \end{equation}
        and then diagonalized.
\end{enumerate}
The maximum eigenfrequency deviation
\begin{equation}
  \max_{n,\vc{q}}
  |\omega_n^\mathrm{ph} - \omega_n^\mathrm{rec}|
\end{equation}
serves as a direct measure of the quality of the extraction.  In
addition, overlay plots of the two band structures provide a visual
consistency check.

% ======================================================================
\section{The \texttt{phonon\_inputs} Package}
\label{sec:phonon_inputs_pkg}
% ======================================================================

The \texttt{phonon\_inputs} package implements the full sow-run-reap
workflow for harmonic and cubic force constants, the structural
relaxation step that precedes it, the structure / configuration
plumbing, and the writers that emit the dynamical-matrix blocks
consumed by the downstream NEGF solver.  The package is calculator-
agnostic: any module that touches DFT goes through a thin adapter
layer that supports Quantum~ESPRESSO \texttt{pw.x} and VASP equally,
and the FC-generation backends (phono3py + \texttt{symfc}, DFPT via
\texttt{ph.x} + \texttt{d3q.x}, and hiphive) share the same dataclass-
based YAML schema and the same on-disk metadata format.

This section walks through the package module by module.  The intent
is to give the reader a single place to look up which module owns
which responsibility, what data flows into it, and what it returns to
the caller.  Module names are given relative to
\path{phonon/phonon_inputs/}.

% ----------------------------------------------------------------------
\subsection{Top-level orchestration: \texttt{cli.py}, \texttt{pipeline.py}}
\label{sec:pinputs_cli}
% ----------------------------------------------------------------------

\paragraph{\texttt{cli.py}.}  An \texttt{argparse}-based command-line
front-end exposed via \texttt{python -m phonon\_inputs <command>}.  The
commands implemented are:

\begin{itemize}
  \item \texttt{generate} -- legacy ``one-shot'' flow that builds a
        phonopy supercell, runs all required QE displacements, and
        produces the quatrex input files in one go (kept as a smoke
        test for tightly coupled bulk runs).
  \item \texttt{extract-blocks} -- reads an existing phonopy reap and
        re-extracts the Convention~B blocks (cf.\
        \cref{sec:convention}) without re-running DFT.
  \item \texttt{validate} -- runs the harmonic-model validators of
        \texttt{validation.py} on an existing reap.
  \item \texttt{fc3-sow}, \texttt{fc3-run}, \texttt{fc3-reap},
        \texttt{fc3-all} -- the canonical phono3py + \texttt{symfc}
        flow.
  \item \texttt{fc3-hiphive-sow}, \texttt{fc3-hiphive-run},
        \texttt{fc3-hiphive-reap},
        \texttt{fc3-hiphive-bootstrap-reap},
        \texttt{fc3-hiphive-convergence},
        \texttt{fc3-hiphive-all} -- the hiphive flow, with the
        two-stage bootstrap variant for phonon-mode rattling.
  \item \texttt{dfpt-sow}, \texttt{dfpt-run}, \texttt{dfpt-reap},
        \texttt{dfpt-all} -- DFPT via \texttt{ph.x} + D3Q.
  \item \texttt{pipeline} -- relax + FC generation in a single
        invocation; calls the in-process relax driver from
        \texttt{qe\_interface.py} and threads the relaxed cell directly
        into the chosen FC backend.
\end{itemize}

All sow-type commands accept the flag \texttt{--from-yaml}.  Without
that flag the structure used for sowing is loaded from the relaxed
output of \texttt{config.relax.work\_dir} (\texttt{CONTCAR} for VASP,
\texttt{relax.out} for QE) via
\texttt{structure.resolve\_structure\_for\_sow}; with the flag, sowing
falls back to the YAML structure as-is.  The default fails fast with a
clear error message when the relax dir is missing, removing a class
of silent failures where a sow ran on the unrelaxed YAML positions and
the resulting FC2 absorbed residual linear forces as artificial
negative curvature.

\paragraph{\texttt{pipeline.py}.}  Implements the \texttt{pipeline}
subcommand.  Two stages: (i) call \texttt{qe\_interface.run\_vasp\_relax}
or \texttt{run\_qe\_relax} on the YAML structure and obtain a relaxed
\texttt{PhonopyAtoms}; (ii) dispatch to the FC backend selected by
\texttt{config.fc\_method} (``finite\_displacement'', ``dfpt'', or
``hiphive'').  Because the relaxed cell is passed in memory directly to
the sow function the pipeline path does not need to write a CONTCAR to
disk to communicate the equilibrium geometry.

\paragraph{\texttt{\_\_main\_\_.py}.}  Six-line shim:
\texttt{from .cli import main; main()}.  Enables
\texttt{python -m phonon\_inputs} as the entry point.

% ----------------------------------------------------------------------
\subsection{Configuration: \texttt{config.py}, \texttt{constants.py}}
\label{sec:pinputs_config}
% ----------------------------------------------------------------------

\paragraph{\texttt{config.py}.}  A nested dataclass schema parsed from
YAML by \texttt{load\_config}.  The top-level class
\texttt{PhononInputConfig} aggregates:

\begin{description}
  \item[\texttt{StructureConfig}.]  \texttt{source}
        (\texttt{inline}\,/\,\texttt{cif}\,/\,\texttt{poscar}\,/\,\texttt{phonopy\_yaml}\,/\,\texttt{qe\_input}),
        with the inline branch providing \texttt{symbols},
        \texttt{lattice}, and \texttt{scaled\_positions} directly.
  \item[\texttt{SupercellConfig}.]  \texttt{matrix},
        \texttt{displacement\_distance} for the phonopy finite-
        displacement backend.
  \item[\texttt{QEConfig} / \texttt{VASPConfig}.]  All knobs the DFT
        backend writes into the input file: \texttt{ecutwfc},
        \texttt{ecutrho\_factor}, \texttt{kpoints\_scf}, \texttt{conv\_thr},
        \texttt{smearing}, \texttt{degauss}, \texttt{pseudo\_dir},
        \texttt{pseudopotentials}, and the analogous VASP fields
        \texttt{encut}, \texttt{ediff}, \texttt{ismear}, \texttt{sigma},
        \texttt{prec}, \texttt{lreal}, \texttt{ncore}, \texttt{kpar},
        \texttt{potcar\_dir}, \texttt{potcar\_map}.
  \item[\texttt{ThirdOrderConfig}.]  Supercell multipliers,
        \texttt{cutoff\_pair\_distance}, \texttt{displacement\_distance},
        \texttt{fc\_calculator} (\texttt{symfc} by default),
        \texttt{work\_dir}, \texttt{pw\_timeout}.
  \item[\texttt{HiphiveConfig}.]  The full hiphive recipe:
        \texttt{supercell}, \texttt{n\_structures}, \texttt{rattle\_method}
        (\texttt{mc}/\texttt{normal}/\texttt{phonon}), \texttt{rattle\_std},
        \texttt{rattle\_d\_min}, \texttt{rattle\_n\_iter},
        \texttt{rattle\_seed}, \texttt{cutoffs}, \texttt{fit\_method},
        \texttt{fit\_kwargs}, \texttt{rotational\_sum\_rule}, the optional
        \texttt{convergence:} sub-block, and the phonon-rattle bootstrap
        knobs \texttt{phonon\_rattle\_temperature\_k},
        \texttt{phonon\_rattle\_bootstrap\_n},
        \texttt{phonon\_rattle\_seed\_fc3} (use existing
        \texttt{fc3.hdf5} as FC2 seed) and
        \texttt{phonon\_rattle\_seed\_dft\_dir} (refit FC2 seed from an
        external mc-rattle work dir's DFT outputs).
  \item[\texttt{DFPTConfig}.]  \texttt{q\_mesh}, \texttt{kpoints},
        \texttt{tr2\_ph}, and the binary paths for \texttt{ph.x},
        \texttt{q2r.x}, and \texttt{d3q.x}.
  \item[\texttt{RelaxConfig}.]  \texttt{calculation}
        (\texttt{vc-relax}/\texttt{relax}), \texttt{forc\_conv\_thr},
        \texttt{press\_conv\_thr}, \texttt{calculator},
        \texttt{work\_dir}.
  \item[\texttt{BlockExtractionConfig},
        \texttt{QuatrexOutputConfig}.]  Block extraction parameters
        (transport axis, $q$-mesh, amplitude cutoff) and quatrex
        writer parameters (file paths, $k$-point grid).
\end{description}

The loader \texttt{config\_from\_dict} traverses the schema, accepting
nested optional blocks (``\texttt{NestedConfig | None}'').  Unknown
keys raise a clear error rather than being silently dropped.

\paragraph{\texttt{constants.py}.}  All unit conversions and physical
constants in a single place.  The constant
\texttt{CONVERSION\_THZ2}~$= e/(m_u\,\text{\AA}^2)/(2\pi)^2\cdot 10^{-24}$ converts
eV/(\AA$^2\cdot$amu) dynamical-matrix eigenvalues to THz$^2$ -- this
is the canonical internal unit-system used throughout the package and
the solver (\cref{sec:phonon_solver_pkg}).  A legacy \texttt{CONVERSION}
constant produces (rad/s)$^2$ and is preserved for tests that compare
to the previous unit convention.  Boltzmann's constant, $\hbar$, the
electron mass and atomic mass unit appear here in SI.

% ----------------------------------------------------------------------
\subsection{Structure loading: \texttt{structure.py}}
\label{sec:pinputs_structure}
% ----------------------------------------------------------------------

\texttt{structure.py} provides two functions and one helper:

\begin{description}
  \item[\texttt{load\_structure(cfg)}.]  Resolves the YAML
        \texttt{structure:} block into a \texttt{PhonopyAtoms} object.
        Inline structures are passed straight through; file-based
        sources (CIF, POSCAR, phonopy YAML, QE input) are read via the
        relevant phonopy / ASE interface.
  \item[\texttt{resolve\_structure\_for\_sow(cfg, config\_dir, \\
        from\_yaml=False)}.]  The default routing layer for sow-time
        entry points.  When \texttt{from\_yaml=False} (the default) it
        looks for \texttt{config.relax.work\_dir/CONTCAR} (VASP) or
        \texttt{relax.out} (QE) and parses the relaxed structure from
        there; if neither file is present it raises \texttt{SystemExit}
        with a precise message directing the user to either run the
        relax first or pass \texttt{--from-yaml}.  With
        \texttt{from\_yaml=True} the YAML structure is used unchanged.
  \item[\texttt{create\_phonopy(cell, supercell\_matrix,\\
        primitive\_matrix, displacement\_distance)}.]  Builds a
        \texttt{Phonopy} object with the chosen supercell and generates
        symmetry-inequivalent displacements at the given amplitude.
\end{description}

The \texttt{Bohr/Rydberg} -- \texttt{Angstrom/eV} conversion factors
for QE-internal storage are pinned here as well, both to honour
phonopy's QE units (Ry/bohr$^2$) on output and to convert forces
parsed from a QE log to eV/\AA{}.

% ----------------------------------------------------------------------
\subsection{Convention transformation and block extraction: \texttt{convention.py}}
\label{sec:convention}
% ----------------------------------------------------------------------

Phonopy stores the dynamical matrix in what we refer to as
\emph{Convention~A}, with phase factor
$\exp[2\pi i\,\vc{q}\cdot(\vc{R} + \boldsymbol{\tau}' - \boldsymbol{\tau})]$
that includes the basis-atom shifts.  The quatrex NEGF solver expects
\emph{Convention~B}, with phase $\exp[2\pi i\,\vc{q}\cdot\vc{R}]$ and
$D_B(\vc{q}) = P(\vc{q})\,D_A(\vc{q})\,P(\vc{q})^\dagger$ where
$P(\vc{q}) = \mathrm{diag}\!\bigl[\exp(2\pi i\,\vc{q}\cdot\boldsymbol{\tau}_\kappa)\bigr]$
repeated three times per atom.  The inverse discrete Fourier
transform of $D_B(\vc{q})$ produces real-valued real-space blocks
$H(\vc{R})$ that the solver tiles into the device block-tridiagonal.

\texttt{convention.py} implements both transformations
(\texttt{apply\_phase\_correction}, \texttt{remove\_phase\_correction})
and the block-extraction driver \texttt{extract\_blocks(phonon,
q\_mesh, amplitude\_cutoff)}, which iterates over a $q$-grid sampling
the transverse Brillouin zone, computes $D_B(\vc{q})$ from the
phonopy-internal $D_A$, and IDFTs the result onto a real-space
$(R_x, R_y, R_z)$ block dictionary.  Blocks with Frobenius norm
below \texttt{amplitude\_cutoff} times the maximum block norm are
dropped.

% ----------------------------------------------------------------------
\subsection{Force-constant production: \texttt{force\_constants.py}}
\label{sec:fc_production}
% ----------------------------------------------------------------------

\texttt{force\_constants.py} owns the harmonic FC2 reconstruction
step: given a \texttt{Phonopy} object and the parsed forces from the
displaced-supercell DFT runs, it calls
\texttt{Phonopy.produce\_force\_constants} with the
``\texttt{symmetry}'' option enabled to symmetry-adapt the FCs, and
honours the \texttt{produce\_fc} hook from \texttt{phono3py} when an
FC3 model is present.  The output is stored on the
\texttt{Phonopy.force\_constants} attribute, ready for q-point
diagonalisation by the validation routines.

% ----------------------------------------------------------------------
\subsection{DFT interface: \texttt{qe\_interface.py}}
\label{sec:dft_interface}
% ----------------------------------------------------------------------

\texttt{qe\_interface.py} unifies the QE and VASP entry points
behind a small set of helpers.  The QE side covers
\texttt{write\_qe\_scf\_input} (builds a single pw.x SCF or relax
input, encoding cell, atomic species, pseudopotentials and k-points;
also writes restart-data block headers so successive displacements can
reuse prior WFC/CHG when convenient),
\texttt{parse\_qe\_force\_output} (reads forces from a finished pw.x
run, converting Ry/bohr to eV/\AA),
\texttt{parse\_qe\_relax\_output} (parses ``Begin final coordinates''
or fixed-cell final positions blocks into a \texttt{PhonopyAtoms}),
and \texttt{run\_qe\_relax}.  The VASP side covers
\texttt{parse\_vasp\_relax\_output} (reads \texttt{CONTCAR}) and
\texttt{run\_vasp\_relax} (writes INCAR/KPOINTS/POSCAR/POTCAR with
the user's relaxation tolerance, runs VASP, and parses the resulting
\texttt{CONTCAR}).  Lower-level helpers (\texttt{\_run\_and\_tee},
\texttt{\_needs\_shell}) propagate stdout/stderr through Python and
detect whether the VASP command requires a shell wrapper.

% ----------------------------------------------------------------------
\subsection{Phono3py + \texttt{symfc} flow: \texttt{thirdorder.py}}
\label{sec:thirdorder}
% ----------------------------------------------------------------------

\texttt{thirdorder.py} drives the finite-displacement FC2+FC3
pipeline.  Three calculator-agnostic public entry points,
\texttt{sow / run\_displacements / reap}, plus a unifying
\texttt{generate\_fc3}.

\paragraph{Sow.}  Builds a \texttt{Phono3py} object with the requested
supercell and pair cutoff, calls its built-in displacement generator
(canonical symmetry-inequivalent finite displacements), and writes one
DFT input per displacement: pw.x input files for QE or
INCAR/KPOINTS/POSCAR/POTCAR directories for VASP.  A reference
(undisplaced) calculation is sown alongside, used as the WFC/CHG
seed for the displaced calculations to accelerate SCF convergence.
Metadata (number of displacements, supercell layout, displacement
distance, calculator) is written to \texttt{hiphive\_meta.json} in
the work directory.

\paragraph{Run.}  Iterates over the displaced calculations, executes
the DFT command (in-process or via a job-submission wrapper provided
by the user), checks for completion markers (``JOB DONE'' for QE,
the OSZICAR final step for VASP), and skips already-completed
calculations.  A timeout mechanism kills hung jobs.

\paragraph{Reap.}  Parses forces from every displaced calculation,
calls \texttt{phono3py.run\_force\_constants} with the
\texttt{fc\_calculator='symfc'} backend (symmetry-adapted least
squares; \cite{seko_projector-based_2024} for the underlying algorithm), and
extracts FC2 and FC3 in the phono3py compact representation.  The
final tensors are written to \texttt{fc3.hdf5} with datasets
\texttt{fc2} of shape $(N, N, 3, 3)$ in eV/\AA$^2$ and \texttt{fc3} of
shape either $(N_\mathrm{prim}, N, N, 3, 3, 3)$ (compact) or
$(N, N, N, 3, 3, 3)$ (full) in eV/\AA$^3$, depending on the symfc
output mode.
Independently of the atomic real-space FC3 (pair-)cutoff set here, the
downstream device solver imposes a further \emph{nearest-neighbour block}
truncation when it ingests this \texttt{fc3.hdf5}: the loader
(\path{quatrex/phonon/fc3_loader.py}, \texttt{load\_device\_fc3} /
\texttt{fc3\_to\_phi\_blocks}) keeps only block-triplets $(I,J,K)$ of
transport cells with $|I-J|,\,|I-K|,\,|J-K| \leq 1$
(\texttt{nn\_only=True} by default), discarding all longer-range block
triplets.  This NN block truncation is applied even on the stored
block-sparse path (not only on the dense fallback), so the cubic vertex
effectively sees at most one neighbouring transport cell per leg
regardless of the atomic cutoff; the loader warns when the dropped
Frobenius norm exceeds a set threshold.

The VASP-side helpers \texttt{\_write\_vasp\_inputs} and
\texttt{\_refresh\_vasp\_runtime\_inputs} emit ISYM=0 in every
displaced INCAR (rattle breaks the parent symmetry; force
symmetrisation would otherwise produce inconsistent forces), and
write KPOINTS at the user-configured density.

% ----------------------------------------------------------------------
\subsection{DFPT flow: \texttt{dfpt.py}}
\label{sec:dfpt_pkg}
% ----------------------------------------------------------------------

\texttt{dfpt.py} drives the DFPT-based alternative FC pipeline using
\texttt{ph.x} + \texttt{q2r.x} for FC2 and \texttt{d3q.x} +
\texttt{d3\_qq2rr.x} for FC3 \cite{paulatto_anharmonic_2013,fugallo_ab_2013}.  The sow stage
writes a sequence of paired SCF + phonon inputs at every $q$-point of
the configured mesh, plus the q2r control file and the third-order
\texttt{d3q.x} inputs.  The run stage executes them in dependency
order (SCF $\to$ ph $\to$ q2r $\to$ d3q $\to$ qq2rr).  The reap stage
parses the produced real-space FC files into the phono3py-compatible
\texttt{fc3.hdf5} format so downstream consumers do not need to branch
on the FC backend.

% ----------------------------------------------------------------------
\subsection{Hiphive flow: \texttt{hiphive\_fc3.py}, \texttt{hiphive\_convergence.py}}
\label{sec:hiphive_flow}
% ----------------------------------------------------------------------

\texttt{hiphive\_fc3.py} implements rattled-supercell sampling and the
fit-and-project step at the heart of the hiphive workflow
\cite{eriksson_hiphive_2019}.  Public API:

\begin{description}
  \item[\texttt{sow}.]  Generates rattled supercells.  For the default
        \texttt{mc} rattle method, atomic displacements are drawn from
        a Monte-Carlo walk with per-step Gaussian increments of standard
        deviation \texttt{rattle\_std} and a hard
        minimum-pair-distance constraint at \texttt{rattle\_d\_min},
        for \texttt{rattle\_n\_iter} iterations; the cumulative
        per-atom displacement scales as $\sqrt{\texttt{rattle\_n\_iter}}$
        (each MC step adds an independent zero-mean Gaussian increment,
        so the per-step variances add).  For
        \texttt{phonon} rattling, a seed FC2 is required: it can come
        from an existing \texttt{fc3.hdf5}
        (\texttt{phonon\_rattle\_seed\_fc3}), an external mc-rattle work
        directory (\texttt{phonon\_rattle\_seed\_dft\_dir}, which
        triggers an FC2-only refit on those forces), or a local
        ``bootstrap'' stage that runs a small mc-rattle pool first and
        fits FC2 from it.  The seed FC2 is checked for PSD-ness before
        being used to drive the phonon-rattle amplitude (which scales as
        $1/\sqrt{|\omega|}$ per mode and diverges for soft / imaginary
        modes; see \cref{sec:hiphive_pathologies}).  Both rattle
        methods write a calculator-specific DFT input per
        \texttt{disp-XXXXX/} directory and a reference (undisplaced)
        directory.

  \item[\texttt{bootstrap\_reap}.]  Stage~1 of the phonon-rattle
        workflow.  Reads the mc-rattle bootstrap pool's forces, fits
        FC2-only with cutoffs $[r_c^{(2)}]$ using ARDR (with ridge as
        fallback), applies the post-fit rotational projection, checks
        that the resulting FC2 has at most
        \texttt{phonon\_rattle\_max\_seed\_imag} imaginary modes, and
        saves \texttt{fc2\_seed.npy} for stage~2.

  \item[\texttt{reap}.]  Reads forces from every \texttt{disp-XXXXX/}
        directory, builds a \texttt{ClusterSpace} with the configured
        cutoffs, runs the selected fit method via
        \texttt{trainstation.Optimizer}, optionally projects onto the
        rotational manifold
        \texttt{enforce\_rotational\_sum\_rules(cs, params,
        sum\_rules=['Huang', 'Born-Huang'])}, builds the
        \texttt{ForceConstantPotential} and exports FC2 and FC3 to
        \texttt{fc3.hdf5} in the same on-disk format as the phono3py
        reap.  A diagnostic JSON \texttt{hiphive\_fit.json} records
        $n_\mathrm{parameters}$, RMSE train/test, rotational residual
        before/after, and the FC2/FC3 ASR residuals.

  \item[\texttt{generate\_fc3}.]  Orchestrates the full hiphive flow
        for the \texttt{phonon} rattle method (two-stage
        bootstrap-aware sequence) and the \texttt{mc}/\texttt{normal}
        rattle methods (one-stage canonical flow).
\end{description}

The helper \texttt{\_seed\_from\_existing\_fc3} extracts FC2 from an
existing \texttt{fc3.hdf5} (e.g.\ a converged mc-rattle reap on the
same primitive) and writes it as \texttt{fc2\_seed.npy}, allowing the
user to bypass the bootstrap stage entirely with the
\texttt{phonon\_rattle\_seed\_fc3} config option.  The helper
\texttt{\_fit\_fc2\_seed\_from\_source} is the corresponding entry
point for refitting FC2 on an external mc-rattle work directory's
\texttt{disp-XXXXX/} outputs, used by
\texttt{phonon\_rattle\_seed\_dft\_dir}.

\texttt{hiphive\_convergence.py} drives the convergence sweep
described in \cref{sec:underdet}: for each pair of (size, fit method)
in the user-configured grid, it subsamples the master DFT pool,
refits, and records the same diagnostic block as the main reap.  The
post-fit rotational projection is exposed by
\texttt{\_apply\_rotational\_sum\_rules} and
\texttt{\_rotational\_residual\_pair}; the latter returns
$\|M_\mathrm{rot}\theta\|$ before and after projection for the
``\texttt{post\_fit}'' policy.

% ----------------------------------------------------------------------
\subsection{Validation: \texttt{validation.py}}
\label{sec:validation}
% ----------------------------------------------------------------------

\texttt{validation.py} provides the four extraction-time checks
detailed in \cref{sec:extraction_validation}: the $\Gamma$-point ASR
residual, the block-symmetry check $H(\vc{R})^T = H(-\vc{R})$, the
band-structure overlay versus the phonopy reference, and a reference
Sancho--Rubio + Caroli ballistic transmission against which the
quatrex solver is benchmarked.  These run in the units chosen by the
caller; the default is THz$^2$ (\texttt{conversion\_factor=CONVERSION\_THZ2}).

% ----------------------------------------------------------------------
\subsection{Tensor decompositions: \texttt{fc3\_compression.py}, \texttt{separable.py}, \texttt{pcp.py}}
\label{sec:fc3_compression}
% ----------------------------------------------------------------------

\texttt{fc3\_compression.py} houses the six tensor-decomposition
ansätze used to compress the supercell FC3.  All operate on the
mass-weighted target $T \in \mathbb{R}^{n_\mathrm{dof}\times \dim_{sc}\times \dim_{sc}}$
constructed by \texttt{build\_fc3\_target} from the raw
$(n_{\mathrm{prim}}, n, n, 3, 3, 3)$ or $(n, n, n, 3, 3, 3)$ FC3
together with the supercell masses:

\begin{description}
  \item[\texttt{mSVD}.]  Truncated matricisation SVD (Tucker mode-3) of
        $T$ flattened along the third axis.
  \item[\texttt{HOSVD}.]  Higher-order SVD (Tucker decomposition) with
        S$_2$ symmetry on legs 2 and 3, refined by a few steps of
        symmetric HOOI.
  \item[\texttt{CP} / \texttt{INDSCAL} / \texttt{Waring}.]
        Canonical polyadic decompositions of increasing symmetry: CP
        unrestricted, INDSCAL with $b_r = c_r$ (S$_2$ symmetry on legs
        2,~3), Waring symmetric CP on the fully S$_3$-symmetrised
        lifted target.
  \item[\texttt{PCP}.]  The permanent-CP ansatz of
        \cite{luo_tensor_2025}, wrapped
        in \texttt{pcp.py}.  Decomposes the cluster representation
        directly as
        $\tilde\Phi_{l_1 b_1, l_2 b_2, l_3 b_3}^{a_1 a_2 a_3}
         = \sum_\xi (\lambda_\xi/3!) \sum_{\sigma\in S_3}\sum_l
           \prod_i A^\xi_{\sigma_i}(l_i - l, b_i, a_i)$
        with three mode functions per rank.
\end{description}

The helper \texttt{asr\_residual(T\_hat, n\_super)} computes the axis-$j$
and axis-$k$ Frobenius residuals
$\|\sum_j T_{ijk}\|_F$ and $\|\sum_k T_{ijk}\|_F$ of the lifted
tensor; \texttt{asr\_project\_factor} performs the corresponding
projection onto the axis-$j$ + axis-$k$ null-space when the downstream
consumer (e.g.\ the bubble in \cref{sec:phonon_solver_pkg}) needs
strict ASR on those legs.

\texttt{separable.py} implements the stacked-SVD low-rank
factorisation of FC3 used by the q-resolved solver
(\cref{sec:se_q}): for each primitive-cell row $(i, \mu)$ it
flattens the mass-weighted slice into a matrix
$M_{(i,\mu)} \in \mathbb{R}^{n\cdot 3 \times n \cdot 3}$, stacks the
slices into $M_\mathrm{stacked}$, factors $M_\mathrm{stacked} =
U_R\,\mathrm{diag}(\sigma)\,V_R^T$ by truncated SVD, and returns
the rank-$R$ left/right factors.

% ----------------------------------------------------------------------
\subsection{Quatrex output: \texttt{quatrex\_writer.py}}
\label{sec:quatrex_writer}
% ----------------------------------------------------------------------

\texttt{quatrex\_writer.py} emits the three files the production
quatrex NEGF stack consumes: \texttt{dynamical\_matrix.mat} (a sparse
MATLAB-format block dictionary keyed by $(R_x, R_y, R_z)$,
\texttt{structure.xyz} (atomic coordinates in extended XYZ for input
to the simulation), and \texttt{config.toml} (the simulation-side
configuration mirror).  All paths and the transport direction are read
from \texttt{QuatrexOutputConfig}.

% ----------------------------------------------------------------------
\subsection{I/O and bookkeeping: \texttt{io.py}}
\label{sec:io_helpers}
% ----------------------------------------------------------------------

\texttt{io.py} provides a small NPZ + JSON sidecar pattern used by the
SCBA and decomposition drivers: arrays go to \texttt{.npz}, scalars and
metadata to a sibling \texttt{.json}.  \texttt{save\_transport\_results}
and \texttt{load\_transport\_results} round-trip the SCBA output
dictionary; an additional helper writes the per-decomposition
``cache'' entries that \texttt{transport\_quality} in
\texttt{finite\_analysis} (\cref{ap:finite_analysis}) reads on
restart.


% ======================================================================
\section{Imaginary-mode pathologies in hiphive SiNW fits}
\label{sec:hiphive_pathologies}
% ======================================================================

This section documents a systematic audit of the imaginary-mode
problem observed for the H-passivated Si nanowire (SiNW) family used
throughout the rest of this report.  Even
after a sequence of cleanups to the wire generator and the relax/sow
plumbing (tighter \texttt{EDIFF}, $a_\mathrm{lattice}$ set to the
PBE-equilibrium value, auto-sized vacuum, FC2 cutoff scaled with the
wire envelope), the regenerated \texttt{cluster/reaps/} still showed
imaginary modes in the resulting harmonic spectrum at $\Gamma$ and at
non-$\Gamma$ wave\-vectors.  We localised the breakage to four distinct
mechanisms, each acting on a different system and each requiring a
targeted fix.  The raw evidence (JSON artefacts, refit sweeps,
species-resolved eigenvalues) lives in
\path{phonon/scratch/imag_audit/}.

% ----------------------------------------------------------------------
\subsection{What ``negative squared frequency'' means physically}
\label{sec:imag_phys}
% ----------------------------------------------------------------------

A harmonic mode has frequency $\omega^2 = \lambda$ where $\lambda$ is
the eigenvalue of the mass-weighted dynamical matrix
$D = M^{-1/2}\,\Phi_2\,M^{-1/2}$.  By convention dispersion plots
report $\mathrm{sign}(\lambda)\sqrt{|\lambda|}$, so a "negative
frequency" is shorthand for $\lambda < 0$ -- the mode is at a saddle
or maximum of the Born-Oppenheimer surface along the corresponding
eigenvector.  The equation of motion is then $\ddot{x} = -\lambda x$,
which gives exponential runaway $x(t) \sim e^{\sqrt{|\lambda|}\,t}$
along that mode coordinate.

Imaginary modes are physical in soft-mode systems (PbTiO$_3$ above its
ferroelectric transition, anharmonically-stabilised bcc-Ti at high
temperature, charge-density-wave instabilities)
\cite{tadano_first-principles_2018}, but for an H-passivated SiNW at its ground-state
geometry they should not appear in the harmonic spectrum.
\Cref{tab:imag_mode_triage} lays out the triage we use to discriminate
real soft-mode physics from numerical artifacts in our setting.

\begin{table}[ht]
  \centering
  \begin{tabular}{lll}
    \toprule
    Magnitude $|\omega|$ & Location & Reading \\
    \midrule
    $< 0.01$ THz & $\Gamma$, three modes & ASR roundoff, ignore \\
    $0.01$--$0.5$ THz & near-$\Gamma$ acoustic & sampling noise on a soft branch \\
    $0.5$--$5$ THz & finite $q$, few modes & cutoff or under-determination \\
    $> 5$ THz & many modes & ground-truth or wiring bug \\
    \bottomrule
  \end{tabular}
  \caption{Imaginary-mode triage matrix.  Magnitudes are mass-weighted
  $\omega$ on the supercell q-mesh.  The $\Gamma$-point three-acoustic-mode
  ASR-noise band is set by the numerical conditioning of the FC fit;
  see \cref{sec:rattle_amplitude}.}
  \label{tab:imag_mode_triage}
\end{table}

% ----------------------------------------------------------------------
\subsection{The four mechanisms}
\label{sec:four_mechanisms}
% ----------------------------------------------------------------------

\Cref{tab:imag_mechanisms} summarises the four root causes we
identified, the system on which each one is dominant, and the
intervention that fixes it.  The rest of this section justifies each
entry with quantitative evidence.

\begin{table}[ht]
  \small
  \centering
  \begin{tabular}{p{0.16\linewidth} p{0.42\linewidth} p{0.34\linewidth}}
    \toprule
    Mechanism & Symptom & Fix \\
    \midrule
    FC2 cutoff $<$ wire diameter
      & d5a sc4: 7 $\Gamma$-imag modes ($\omega \in [-2.2,-1.2]$ THz), 90 \%
        Si-localised.
      & Set cutoffs[0] $\geq 2 \times \mathrm{envelope} + 0.5$ Å. \\
    \addlinespace
    Under-determination at large cutoffs
      & d9a sc4 at cutoffs=[7.83, 4.0]: $n_\mathrm{eqn}/n_\mathrm{params} = 0.70$
        (40 \% under-det).  Validator already errors.
      & Bump $n_\mathrm{structures}$ to $4 \lceil n_\mathrm{super}/8 \rceil$.
        Re-DFT required. \\
    \addlinespace
    Over-aggressive mc-rattle on light H atoms
      & d9a sc4 with $n_\mathrm{iter}=10$: rattle peaks (max over atoms)
        0.27–0.37 Å,
        Si–H bonds compressed 13 \%, anharmonic forces absorbed as
        artificial negative FC2 curvature on the H subspace
        (51 $\Gamma$-imag modes at –14 THz, 91 \% H-localised).
      & $n_\mathrm{iter}=2$, std=0.03 Å (cumulative $\approx 0.06$ Å). \\
    \addlinespace
    Stale YAML positions vs the relaxed sow geometry
      & 0.55 Å drift between \texttt{hiphive\_meta.json} (sow time) and
        the \texttt{configs/sinw/*.yaml} the analysis pipeline reads.
        Affects $q \neq \Gamma$ phase factors.
      & In \texttt{loader.py}, source the primitive from
        \texttt{hiphive\_meta.json} when present. \\
    \bottomrule
  \end{tabular}
  \caption{The four mechanisms we identified and the fixes applied in
  the code review.  ``d5a'' and ``d9a'' refer to the
  Si$_9$H$_{12}$ and Si$_{21}$H$_{16}$ wires in
  \texttt{phonon/configs/sinw/}.  Atom counts, lattice
  parameters and cross-section conventions follow
  \cite{Rurali2010colloquium,Peelaers2009sinw}.}
  \label{tab:imag_mechanisms}
\end{table}

% ----------------------------------------------------------------------
\subsubsection{Mechanism 1: FC2 cutoff vs cross-section span}
\label{sec:cutoff_vs_section}
% ----------------------------------------------------------------------

A hiphive ClusterSpace built with $\mathrm{cutoffs}[0] = r_c^{(2)}$
assigns zero FC2 to any atom pair separated by more than $r_c^{(2)}$ in
Cartesian distance \cite{eriksson_hiphive_2019}.  For a wire of
H-shell radius $\rho$, the cross-section spans up to $2\rho$
(top-to-bottom of the H envelope), so a cutoff with
$r_c^{(2)} \lesssim 2\rho$ leaves transverse Si--Si pairs unconstrained.
The corresponding modes -- breathing, ovalisation, transverse acoustic
-- have no restoring force in the fit basis and either localise (no
dispersion) or pick up a finite negative dynamical-matrix eigenvalue
through coupling with rotationally-projected FC2 entries.

A cutoff sweep on the d5a sc4 cluster reap (84 atoms, fixed at
$n_\mathrm{structures}=16$, refit on the existing DFT outputs)
quantifies this.  Each cell of \cref{tab:d5a_cutoff_sweep} is the
$\Gamma$-point imaginary-mode count of the resulting
mass-weighted dynamical matrix; ARDR and least-squares fits with and
without the post-fit Huang/Born-Huang projection are reported
separately.  The values come from
\path{phonon/scratch/imag_audit/pass1_cutoff_sweep.json}.

\begin{table}[ht]
  \centering
  \begin{tabular}{lcccc}
    \toprule
    $r_c^{(2)}$ (Å) & ARDR / no RSR & ARDR / RSR & LS / no RSR & LS / RSR \\
    \midrule
    5.0 & 13 & 16 & 16 &  6 \\
    5.5 &  4 & 14 & 16 & 13 \\
    6.0 &  4 &  7 &  8 & 11 \\
    6.5 &  4 &  6 & 13 & 16 \\
    7.0 &  4 & \textbf{0} & 15 & 16 \\
    7.8 &  4 & \textbf{0} & 15 & 18 \\
    \bottomrule
  \end{tabular}
  \caption{$n_\mathrm{imag}$ at $\Gamma$ as a function of FC2 cutoff for
  the d5a sc4 wire (envelope $\rho = 3.7$ Å, full diameter $\approx 7.4$ Å).
  ARDR + post-fit rotational sum rule succeeds only once
  $r_c^{(2)} \geq 7.0$ Å, i.e.\ once the cutoff spans nearly the full
  H-envelope diameter of $\approx 7.4$ Å.  All fits use FC3 cutoff
  4.0 Å, $n_\mathrm{structures}=16$.}
  \label{tab:d5a_cutoff_sweep}
\end{table}

Two observations are robust.  First, ARDR without the rotational
projection picks up a fixed-and-small $n_\mathrm{imag}=4$ for any
$r_c^{(2)} \geq 5.5$ Å -- its sparse selection keeps FC2 close to PSD
but does not enforce rotational invariance.  Second, the Huang +
Born-Huang projection succeeds only once $r_c^{(2)} \geq$ the wire
diameter; at smaller cutoffs the projection enforces rotational
invariance perfectly but shifts FC2 onto a less-PSD manifold and
\emph{introduces} imaginary modes.
% REVIEW(open): bib entry needed for the rotational-sum-rule vs PSD trade-off claim (the original Mounet2005phonons was unresolvable and off-topic; citation dropped)

Refitting the d5a sc4 supercell at $r_c^{(2)} = 7.0$ Å with ARDR + RSR
(the smallest clean value of the sweep) gives a clean dispersion across
the full $\Gamma\to Z$ path: no
imaginary modes across 21 q-points and 63 bands, with
$\min \omega \approx -10^{-5}$ THz (machine epsilon) and
$\max \omega = 62.9$ THz consistent with Si--H stretching.  The
practical consequence for the wire generator is the formula change
implemented in \texttt{setup\_sinw100.py:\_suggest\_fc2\_cutoff},
\begin{equation}\label{eq:suggested_cutoff}
  r_c^{(2)} = \min\bigl(12.0,\ \max(5.0,\ 2\rho + 0.5)\bigr) \quad [\text{\AA}],
\end{equation}
which replaces the previous $\rho + 2.4$ formula that systematically
undershoots.  For the production d5a sc4 configuration this formula
(with the measured envelope $\rho=3.69$ Å) yields the FC2 cutoff
$r_c^{(2)} = 7.88$ Å actually used (\texttt{cutoffs: [7.88, 4.0]} in
\path{phonon/configs/sinw/sinw100_d5a_vasp_sc4.yaml}); the $7.0$ and
$7.8$ Å values above are sweep grid points, not the production cutoff.

% ----------------------------------------------------------------------
\subsubsection{Mechanism 2: Under-determination at fixed sample count}
\label{sec:underdet}
% ----------------------------------------------------------------------

For the d9a wire the FC2 cutoff bump pushes
$n_\mathrm{parameters}$ past
$n_\mathrm{equations} = n_\mathrm{structures} \cdot n_\mathrm{super}\cdot 3$.
At cutoffs $[7.83, 4.0]$ on the d9a sc4 supercell (148 atoms) the
ClusterSpace expands to 10\,116 free parameters, while
$n_\mathrm{eqn} = 16\cdot148\cdot3 = 7\,104$ -- a ratio of $0.70$, i.e.\
40~\% under-determined.  ARDR's automatic relevance determination
\cite{eriksson_hiphive_2019} sparsifies aggressively, but it minimises
the residual training error and cannot guarantee that the FC2
extracted from the sparse parameter subset has a positive-definite
dynamical matrix.  Empirically the result is the d9a sc4 reap's 51
$\Gamma$-imag modes spanning $-14$ to $-2$ THz.

The validator built into \texttt{phonon\_inputs/parameter\_validation}
% REVIEW(verified): the implemented rule target = max(6, 4*(n_super+7)//8) gives 44 for 84 atoms (parameter_validation.py:210), matching eq:nstruct_recommendation; the previously quoted 45 matched no version of the formula in git history
already flags this with severity \texttt{warn} (it raised
``\texttt{n\_structures should be $\geq$ 44 for 84-atom supercell}'' on
the d5a sc4 reap before the present round of fixes).  The recommended
$n_\mathrm{structures}$ scales with the supercell size as
\begin{equation}\label{eq:nstruct_recommendation}
  n_\mathrm{structures}^\mathrm{rec}
  = 4 \cdot \bigl\lceil n_\mathrm{super}/8 \bigr\rceil,
\end{equation}
yielding 44 for d5a sc4, 76 for d9a sc4, and 104 for d12a sc4 -- the
values now hard-coded in
\texttt{configs/sinw/sinw100\_*\_vasp\_sc4*.yaml}.  This is consistent
with the rule of thumb in the hiphive paper that FC fits for
under-determined regimes require ``advanced regularized regression''
together with several-times-over\-determined sampling
\cite{eriksson_hiphive_2019}.

\paragraph{Cutoff alone cannot save d9a.}
We also ran a cutoff sweep for d9a sc4 at fixed
$n_\mathrm{structures}=16$ (FC3 cutoff reduced to 3.0 Å for speed;
\path{phonon/scratch/imag_audit/pass1_d9a_quick.json}).  Even at
$r_c^{(2)}=5.0$ Å, where the ratio recovers to 3.57 ($n_\mathrm{params}=1989$,
$n_\mathrm{eqn}=7104$), the dispersion still has 28 imaginary modes at
$\Gamma$ with $\min\omega = -4.1$ THz, and the situation degrades at
larger cutoffs.  So under-determination is necessary but not
sufficient: the residual imaginary modes are driven by mechanism 3,
the over-aggressive rattle.

% ----------------------------------------------------------------------
\subsubsection{Mechanism 3: Rattle amplitude and the H subspace}
\label{sec:rattle_amplitude}
% ----------------------------------------------------------------------

The hiphive paper explicitly recommends rattle amplitudes
``$\approx 0.01$--$0.05$ Å'' for accurate force-constant extraction
and warns that ``it can become difficult to untangle the force
contribution if higher orders are contributing to the forces.  This
leads to higher-order contributions effectively being included in the
fitted force constants thus reducing their accuracy''
\cite{eriksson_hiphive_2019}.  Our previous default
($\mathrm{rattle\_std}=0.03$ Å, $n_\mathrm{iter}=10$) was inside the
recommended range only at first glance.
In hiphive's mc-rattle the
per-atom cumulative displacement scales as
$\sqrt{n_\mathrm{iter}}$ -- each Monte-Carlo step adds an independent
zero-mean Gaussian increment of standard deviation $\mathrm{rattle\_std}$,
so the per-step variances add -- giving a typical realised displacement
after $n_\mathrm{iter}=10$ of
$\mathrm{rattle\_std}\sqrt{n_\mathrm{iter}} \approx 0.1$ Å, while the
extreme value over all atoms reaches $\approx 0.3$ Å (the measured
$0.27$--$0.37$ Å peaks below).  The practical conclusion is unchanged:
even at the $\sqrt{n_\mathrm{iter}}$ rate the realised amplitudes leave
the harmonic regime.

For three d5a sc4 disp directories we measured atomic-displacement
peaks 0.27--0.37 Å and Si--Si nearest-neighbour distances dropping to
2.04 Å -- a 13~\% bond compression from the equilibrium 2.351 Å, with
forces of 6--7 eV/Å (\path{pass2_3_outcar_force_sanity.json}).  This is
solidly in the anharmonic regime for both Si--Si bonds and the lighter
Si--H bonds; the cubic FC3 term contributes appreciably to the force
along these compressed bonds, but the harmonic fit absorbs the cubic
remainder as artificial negative curvature.  The bias preferentially
affects the H subspace because the mass-weighted dynamical-matrix
eigenvalue $\lambda$ scales as $\Phi_2 / m$: an error
$\delta\Phi_2$ on the H--H block produces an eigenvalue shift
$\delta\lambda = \delta\Phi_2/m_\mathrm{H}$, which is
$m_\mathrm{Si}/m_\mathrm{H} \approx 28$ times more sensitive than the
analogous error on the Si--Si block.

\paragraph{Species-localised pathology pattern.}
Diagonalising the d5a sc4 mass-weighted $\Phi_2$ at $\Gamma$
(\path{pass2_d5a_d9a_gamma_eig.json}) gives seven negative eigenvalues
between $-2.2$ and $-1.2$ THz, with Si occupation 70--91 \% on each
eigenvector.  The same diagonalisation for d9a sc4 gives 51 negative
eigenvalues between $-14$ and $-2$ THz, with H occupation 88--91 \%.
The species-localisation pattern is the clean signature of the
mass-leverage mechanism just described.

The fix is to drop $n_\mathrm{iter}$ to 2 (cumulative displacement
$\approx 0.06$ Å, well within the harmonic regime), keeping
$\mathrm{rattle\_std}=0.03$ Å.  This matches the hiphive
thermal-conductivity tutorial defaults and the displacement-amplitude
guidance in \cite{eriksson_hiphive_2019}.

% ----------------------------------------------------------------------
\subsubsection{Mechanism 4: Stale YAML positions}
\label{sec:stale_yaml}
% ----------------------------------------------------------------------

A subtle plumbing bug: the analysis-time \texttt{Phonopy} object in
\texttt{finite\_analysis/loader.py:\_build\_phonopy} was constructed
from \texttt{cfg.structure}, the YAML's primitive cell, while the
fitted FC arrays in \texttt{fc3.hdf5} were produced from the
\emph{relaxed} primitive sown by
\texttt{resolve\_structure\_for\_sow}.  For the d9a wire the relax
displaced atoms by up to 0.55 Å (\path{pass2_phonopy_vs_hiphive_atoms.json}),
so phonopy was diagonalising the dynamical matrix with the wrong atomic
positions in the q-phase factors $\exp(i\vc{q}\cdot \vc{r}_{j}(\vc{R}))$.

At $\Gamma$ this is irrelevant -- the phases collapse to unity -- so
the $\Gamma$-imag count is unaffected.  But the entire dispersion path
at $\vc{q} \neq 0$ silently used wrong phases, scrambling the $\Gamma$-to-Z
band structure.  The fix is to source the primitive from
\texttt{hiphive\_meta.json} when present:
\texttt{loader.py:\_primitive\_from\_reap\_metadata} now does this, and
\texttt{\_build\_phonopy} accepts a \texttt{primitive\_override} that
\texttt{load\_system} passes in.  Bulk-Si (phono3py) reaps without a
sibling \texttt{hiphive\_meta.json} fall back to the YAML primitive,
preserving backward compatibility.

% ----------------------------------------------------------------------
\subsection{The \texttt{leg\_j\_rel} metric is a hiphive structural artefact}
\label{sec:leg_j_rel}
% ----------------------------------------------------------------------

The reading
\texttt{physical/physical.json["fc3\_asr\_legs"]["leg\_j\_rel"]}
$\approx 0.80$ appears in every SiNW hiphive fit and is structural,
not a fit defect.  The metric measures
\begin{equation}
  \mathrm{leg\_j\_rel} =
  \frac{\bigl\| \sum_j \Phi^3_{ijk,\alpha\beta\gamma} \bigr\|_F}{\| \Phi^3 \|_F},
\end{equation}
i.e.\ the relative Frobenius norm of the axis-$j$ trace of the
mass-weighted lifted FC3 tensor.  Hiphive's \texttt{ClusterSpace}
construction enforces the axis-$i$ ASR by including the orbit-symmetry
projector at basis-construction time, but it does \emph{not} enforce
the analogous constraint along axes $j$ or $k$ -- those would require
either a constrained cluster-space basis or an explicit post-fit
projection on the FC tensor itself.

Bulk-Si reaps with full octahedral point symmetry yield
$\mathrm{leg\_j\_rel} = 0.000$ to floating-point precision
(\path{pass3_bulk_si_leg_j.json}); quasi-1D SiNW reaps with a single
mirror symmetry yield $\approx 0.80$ regardless of $\mathrm{fit\_method}$,
$n_\mathrm{structures}$, or cutoffs.  The metric is therefore a
structural feature of hiphive's cluster-space construction on a
low-symmetry quasi-1D system, not a fit defect.

A practical consequence: the SSE bubble in
\texttt{phonon/finite\_analysis/} integrates FC3 along axis $j$
(self-energy at frequency $\omega$ depends on $\sum_j \Phi^3_{ijk}$
contracted with $G^<(\omega-\omega')$), so a finite
$\mathrm{leg\_j\_rel}$ propagates to a finite Drude-like weight in
$\Sigma_\mathrm{ph\text{-}ph}$ as $\omega \to 0$.  If strict axis-$j/k$
ASR is required downstream, project explicitly via
\texttt{phonon\_inputs.fc3\_compression.asr\_project\_factor} (or pass
\texttt{asr\_project=True} to \texttt{load\_quatrex\_blocks}).  The
\texttt{fc3\_asr\_legs} docstring and \texttt{HIPHIVE\_FITTING\_NOTES}
now reflect this.

% ----------------------------------------------------------------------
\subsection{Post-fit rotational projection is FC2-only in practice}
\label{sec:rsr_fc2_only}
% ----------------------------------------------------------------------

The Huang + Born-Huang post-fit projection
$\Pi: \theta \mapsto \arg\min_{\theta'} \|\theta'-\theta\|^2$ subject
to $M_\mathrm{rot}\theta' = 0$ (with $M_\mathrm{rot}$ the rotational
constraint matrix exposed by
\texttt{hiphive.core.rotational\_constraints.\allowbreak get\_rotational\_constraint\_matrix})
operates on the unified cluster-space parameter vector and rebuilds
both FC2 and FC3 from the projected parameters.  The relevant
constraint structure couples FC2 and FC3 via the Born-Huang relation
through cubic orbits of the cluster space \cite{eriksson_hiphive_2019}.

For the FC3 cutoff of 4 Å used throughout this work, the projection
touches only FC2.  Refitting d5a sc4 at cutoffs $[6.09, 4.0]$ and
measuring before/after Frobenius distances
(\path{pass4_d5a_rsr_delta.json}) gives
\begin{align*}
  \|\Phi_2^\mathrm{after} - \Phi_2^\mathrm{before}\|_F / \|\Phi_2\|_F &= 3.0 \times 10^{-2}, \\
  \|\Phi_3^\mathrm{after} - \Phi_3^\mathrm{before}\|_F / \|\Phi_3\|_F &= 0.0,
\end{align*}
even though the rotational residual $\|M_\mathrm{rot}\theta\|$ drops
from 48.8 to $3 \times 10^{-3}$ -- four orders of magnitude.  The
cubic orbits at our $r_c^{(3)} = 4$ Å apparently do not span the
parameter directions that $M_\mathrm{rot}$ selects, so the projection
acts as $\Pi_2 \oplus \mathbb{1}_3$ in the FC-tensor space.

A side effect is the imaginary-mode trade-off displayed in
\cref{tab:d5a_cutoff_sweep}: at small FC2 cutoffs the projection
enforces rotational invariance perfectly but pushes FC2 off the PSD
cone, \emph{adding} $\Gamma$-imaginary modes.  At cutoffs that span the wire
diameter the projection succeeds without that side-effect, and the
$\Gamma$-imag count drops to zero.

% ----------------------------------------------------------------------
\subsection{Are these wires sensible objects?  Comparison to literature}
\label{sec:sinw_literature}
% ----------------------------------------------------------------------

The Si nanowires we use as test systems for the
phonon--phonon-scattering self-energy machinery are at the small end
of the diameter range studied by first-principles phonon work.  For
context:

\begin{itemize}
  \item Peelaers et al.\ \cite{Peelaers2009sinw} compute ab-initio
        phonon band structures of $\langle 110 \rangle$ H-passivated
        SiNWs with diameters $\sim 1$--$1.5$ nm and report that all
        studied wires are dynamically stable provided the unit cell
        along the wire axis is taken sufficiently large.  Their
        smallest wires are similar in cross-section to our d9a and
        d12a structures.
  % REVIEW(open): confirm Rurali2010colloquium actually enumerates the Si9H12/Si21H20/Si37H28 family (original key CartoixaRurali2010 was unresolvable and remains missing from the bib)
  \item The canonical $\langle 100 \rangle$ H-passivated SiNW
        family is Si$_9$H$_{12}$, Si$_{21}$H$_{20}$, Si$_{37}$H$_{28}$, $\ldots$~\cite{Rurali2010colloquium}.
        Our d5a wire is identical to the smallest member of this
        family; our d9a and d12a are circular carves (Si$_{21}$H$_{16}$
        and Si$_{29}$H$_{22}$) that don't fall on the canonical
        $\{100\}/\{110\}$-faceted square cross-section but use the same
        bond chemistry, the same Si--H equilibrium of 1.48--1.50 Å
        \cite{Vo2008atomistic}, and identical [100] axis orientation.
  \item Vo, Williamson and Galli \cite{Vo2008atomistic} and Donadio and
        Galli \cite{Donadio2010thermal} use SiNW diameters in the
        $1$--$3$ nm range for thermal conductivity studies; the
        cross-section-shape sensitivity they document supports our
        decision to use circular carves alongside the canonical
        faceted family.
  \item Karamitaheri et al.\ \cite{Karamitaheri2013mvff} study SiNW
        phonon dispersions with side lengths 1--10 nm using a modified
        valence force field; they confirm that confinement-induced
        flattening of acoustic branches sets in below $\sim 3$ nm,
        which is the regime our d5a sits in firmly.
  \item Markussen, Jauho and Brandbyge \cite{Markussen2007anisotropic}
        use $\langle 100 \rangle$ H-passivated SiNWs in the
        $\sim 1$--$3$ nm diameter range as test systems for ballistic
        phonon transport -- exactly the use we put these wires to
        downstream.
\end{itemize}

Two qualifications follow.  First, the d5a wire (9 Si in
cross-section, $\sim 0.5$ nm of Si core, $\sim 0.7$ nm with the H
shell) is closer to a Si chain than a 3D wire.  Quantum confinement
on the electrons is of order 1 eV; flexural branches are not
continuum-Euler-Bernoulli but tight-binding-chain-like.  Quantitative
comparison to experimental SiNW transport data (which is at 10--100 nm
diameters) is not the goal; demonstrating the Σ$_\mathrm{ph-ph}$
infrastructure is.  Second, the under-passivation we see at the
corners of d9a and d12a (3-coordinated Si atoms; cf.\
\texttt{phonon/examples/\_nanowire.py:\_check\_coordination}) is
intrinsic to the circular carve.  A canonical Cartoixà--Rurali
Si$_{21}$H$_{20}$ carve would have all surface Si 4-coordinated;
switching to faceted carves is an optional follow-up if we want to
match the canonical convention exactly.

% ----------------------------------------------------------------------
\subsection{Verification protocol after the fixes}
\label{sec:verif_post_fix}
% ----------------------------------------------------------------------

The full audit produced a punch list of seven items (A--G), summarised
in \path{phonon/scratch/imag_audit/INVESTIGATION_NOTES.md}.  All seven
were implemented in the code review documented elsewhere in this
appendix.  The expected verification path is:

\begin{enumerate}
  \item \textbf{Refit d5a sc4 only} (no DFT, the existing
        \texttt{disp-XXXXX/} forces suffice).  After the FC2 cutoff
        bump to 7.88 Å and the rotational projection, the
        $\Gamma\to Z$ dispersion is clean across 21 q-points × 63
        bands with $\min\omega \approx -10^{-5}$ THz.  Already
        verified in \path{pass1_d5a_best_dispersion.json}.
  \item \textbf{Re-sow + re-DFT d9a sc4} with
        $n_\mathrm{structures}=76$, $r_c^{(2)}=11.36$ Å,
        $n_\mathrm{iter}=2$.  Expected: $n_\mathrm{imag}=0$,
        $\max_{\mathrm{acoustic}@\Gamma} < 0.1$ THz.
  \item \textbf{Re-sow + re-DFT d12a sc4} with
        $n_\mathrm{structures}=104$, $r_c^{(2)}=12$ Å (capped),
        $n_\mathrm{iter}=2$.
  \item For all systems: confirm \texttt{leg\_j\_rel} stays at
        $\sim 0.80$ (structural, expected); do not treat as a fit
        defect.
\end{enumerate}

Item 1 is hours of refitting; items 2 and 3 are cluster runs
(76 × 148-atom + 104 × 204-atom DFT calculations).  These are tracked
as TODOs and will be filled into \cref{sec:res_fc} once the
cluster re-runs complete.
